\section{Findings and Exploratory Analysis}
\label{sec:findings}

The findings use only loaded observations and disclose the grain and sample
rules behind each result. They describe the database; they do not by themselves
establish national causation because station coverage changes over time.

\subsection{Temperature Change Across the Loaded Series}
\label{subsec:temperature-findings}

Figure~\ref{fig:annual-temperature} plots only years meeting the stated
minimum observation count.

\begin{figure}[H]
\centering
\begin{tikzpicture}
\begin{axis}[
  width=0.90\textwidth,
  height=6.5cm,
  xlabel={Year},
  ylabel={Mean monthly maximum temperature ($^\circ$C)},
  grid=major,
  xmin=1968,xmax=2025,
  legend style={at={(0.5,-0.18)},anchor=north}
]
\addplot+[cublue,thick,mark=none]
  table[x=year,y=value]{scripts/generated/annual_temperature.dat};
\legend{Years with at least 100 loaded observations}
\end{axis}
\end{tikzpicture}
\caption{Annual mean of loaded monthly maximum-temperature observations.
Source: reproducible SQLite query in \texttt{report/measure.py}.}
\label{fig:annual-temperature}
\end{figure}

Among years with at least 100 maximum-temperature observations, the first five
available annual means average \mTempEarlyMean$^\circ$C and the latest five
average \mTempRecentMean$^\circ$C, a descriptive difference of
\mTempDifference$^\circ$C. This is an observed change in the loaded records,
not a homogeneous national temperature series: station membership and
observation counts differ between years.

\subsection{Rainfall Variation and Station Ranking}
\label{subsec:rainfall-findings}

Figure~\ref{fig:annual-rainfall} shows annual variation while
Table~\ref{tab:rainfall-ranking} adds the sample-size-controlled station
comparison.

\begin{figure}[H]
\centering
\begin{tikzpicture}
\begin{axis}[
  width=0.90\textwidth,
  height=6.5cm,
  xlabel={Year},
  ylabel={Mean loaded monthly rainfall (mm)},
  grid=major,
  xmin=1947,xmax=2025
]
\addplot+[cuteal,thick,mark=none]
  table[x=year,y=value]{scripts/generated/annual_rainfall.dat};
\end{axis}
\end{tikzpicture}
\caption{Annual mean of non-null monthly rainfall observations. Source:
reproducible SQLite query in \texttt{report/measure.py}.}
\label{fig:annual-rainfall}
\end{figure}

Rainfall varies more sharply from year to year than the temperature series. A
station ranking is meaningful only when sample size is included. With a
minimum of 100 observations, \mWettestStation{} ranks first with
\mWettestMeasurements{} measurements, an average of \mWettestAverage{} mm and a
maximum reported month of \mWettestMaximum{} mm.

\begin{table}[H]
\centering
\caption{Highest average monthly rainfall among stations with at least 100
measurements}
\label{tab:rainfall-ranking}
\begin{tabular}{L{0.31\textwidth}r r r}
\toprule
Station & Measurements & Average (mm) & Maximum (mm) \\
\midrule
\inputrows{scripts/generated/rainfall_ranking.tex}
\bottomrule
\end{tabular}
\end{table}

\subsection{Water-Quality Parameter Coverage}
\label{subsec:water-quality-findings}

The water-quality relation contains six measured parameter types.
Figure~\ref{fig:water-quality-parameters} shows how unevenly those parameters
are represented. Dissolved oxygen (DO), pH and biochemical oxygen demand (BOD)
account for most observations; chemical oxygen demand (COD), salinity and
marine-debris records are more limited. The chart measures database coverage,
not the environmental importance of a parameter.

\begin{figure}[H]
\centering
\begin{tikzpicture}
\begin{axis}[
  xbar,
  width=0.86\textwidth,
  height=6.5cm,
  xlabel={Loaded observations},
  symbolic y coords={MarineDebris,Salinity,COD,BOD,pH,DO},
  ytick=data,
  yticklabels={DO,pH,BOD,COD,Salinity,Marine debris},
  xmin=0,
  nodes near coords,
  bar width=8pt,
  fill=cuteal!60,
  draw=cuteal
]
\addplot table[x=records,y=parameter]
  {scripts/generated/water_quality_parameters.dat};
\end{axis}
\end{tikzpicture}
\caption{Loaded water-quality observations by parameter. Source: reproducible
SQLite query in \texttt{report/measure.py}.}
\label{fig:water-quality-parameters}
\end{figure}

\subsection{Connections Between Tables}
\label{subsec:relationship-findings}

The final relationships support questions that the source files cannot answer
without manual reconciliation. All \mWaterQualityConnected{}
\texttt{Water\_Quality} rows join through \texttt{WQ\_Station\_Name} to
\texttt{River\_Station} and from there to a retained river or water-body
reference. Climate tables can be
combined at the shared station--year--month grain through the
\texttt{Monthly\_Climate\_Summary} view. Forest observations connect to both a
district and a valid fiscal-year span.

For the \mCompleteClimateRows{} station-month rows where maximum temperature,
humidity and rainfall are all present, Pearson correlation is
\mTempHumidityCorrelation{} between maximum temperature and humidity and
\mHumidityRainfallCorrelation{} between humidity and rainfall. The latter is a
moderate positive association in the loaded monthly records. Neither value
shows causation and neither should be generalized beyond the complete-case
subset.

\subsection{Missing Data and Coverage Limits}
\label{subsec:missing-findings}

Table~\ref{tab:missing-measurements} distinguishes retained null measurements
from the separate station-geography limitation.

\begin{table}[H]
\centering
\caption{Null measurements retained in selected fact tables}
\label{tab:missing-measurements}
\begin{tabular}{L{0.38\textwidth}r r r}
\toprule
Field & Missing & Total & Missing share \\
\midrule
\inputrows{scripts/generated/missing_measurements.tex}
\bottomrule
\end{tabular}
\end{table}

Missing source measurements are never converted to zero. Missing values remain
\texttt{NULL} where a retained row contains a nullable attribute, while
observations with no usable measurement are not fabricated as database records.
The most important geographic limitation is separate from the measurement table:
\mUnmappedStations{} of 56 stations lack a supported district mapping.
Their climate observations remain queryable by station but district-level
analysis must disclose the reduced coverage.

Industry coverage is also narrow. The approved final model contains one
industry category and \mIndustryRows{} fiscal observations. This is a finding
about the available source, not evidence that other industries produce no
wastewater. Similarly, each relation's time range in
Figure~\ref{fig:timespans} must be checked before a cross-subject trend is
interpreted.
